% 5.3 — Stratégie de test
\section{Stratégie de test}
\label{sec:tests}

\subsection{Les familles de tests}
\begin{description}[style=nextline,leftmargin=0.6cm,font=\sffamily\bfseries]
  \item[Les exemples de règles] Chacun des \nbexemples{} exemples de la table de règles est
    exécuté par tout le moteur dans une arborescence fixe, et l'impact obtenu est comparé à
    l'impact attendu (\fichier{test\_engine\_rules.py}). Un exemple marqué \cmd{VERIFIED}
    décrit un comportement observé sur des fichiers jetables.
  \item[Les tests de partie] Chaque partie livre son fichier (\fichier{test\_m23a.py},
    \fichier{test\_v17.py}, \fichier{test\_wsec.py}\ldots), écrit pour échouer sur le code
    d'avant.
  \item[Les lignes figées] Pour chaque famille d'outils, les lignes du portail concernées sont
    figées dans \fichier{tests/fixtures/} avec leur décision attendue ; une partie suivante qui
    les fait bouger doit le justifier.
  \item[Les attentes exactes par système] Quand Linux et Windows diffèrent (chemins, filet,
    dialecte), le test l'écrit explicitement, et la restauration est vérifiée octet par octet.
  \item[Les confrontations au réel] Grammaires comparées au \cmd{--help} des outils installés ;
    prédictions SQL comparées à ce que change réellement un serveur PostgreSQL de test ; le
    garde d'appels système lancé pour de vrai.
  \item[Les tests d'inertie] Moteur activé, désactivé, en erreur ou hors budget : les décisions
    du chemin en ombre ne changent pas (horloge simulée pour forcer le dépassement).
\end{description}

\subsection{Un test qui échoue sur le code d'avant}
Un test qui passe aussi sur le code vulnérable ne prouve rien sur la correction. L'exemple
suivant est le test de V19 :

\begin{code}{Le test de V19 (\fichier{tests/test\_m23a.py})}
def test_5_unread_sql_escalates_instead_of_being_allowed(tmp_path):
    proj = _project(tmp_path)
    for cmd in ('psql -c "$(cat q.sql)"', 'mysql -e "$(cat q.sql)"',
                'sqlite3 app.db "$(cat q.sql)"'):
        r = _eval(proj, cmd)
        assert (r.decision.decision, r.recovery_entry) == ("ESCALATE", None), cmd
\end{code}

Pour le vérifier, les tests de M2.3a ont été exécutés contre le commit précédent
(\cmd{99831a9}) : 11 échouent --- ceux qui décrivent les corrections --- et 6 passent : ce
sont les \emph{témoins}, qui vérifient que ce qui était déjà juste le reste.

\begin{terminal}[les tests de M2.3a sur le code d'avant]
$ git archive 99831a9 src | tar -x -C ancien/ && cd ancien
$ python3 -m pytest -q tests/test_m23a.py
FAILED tests/test_m23a.py::test_1_a_runner_goes_to_the_runtime_net
FAILED tests/test_m23a.py::test_5_unread_sql_escalates_instead_of_being_allowed
FAILED tests/test_m23a.py::test_13_file_tools_that_destroy_are_captured
[...]
11 failed, 6 passed in 0.31s
\end{terminal}

\subsection{Les pièges des tests eux-mêmes}
Deux fois, le banc de test du garde du shell a produit des succès \emph{vides} sous Windows :
le chemin du journal d'appels, écrit dans un script bash, contenait des barres obliques
inverses que bash lisait comme des échappements ; les tests « aucun appel ne doit avoir eu
lieu » passaient alors sans rien prouver. Le chemin est désormais dérivé de \cmd{\$0}, ce qui
supprime la classe d'erreur, et un test témoin vérifie que le journal est bien lu. Un test vert
n'est une preuve que si l'on a vu qu'il pouvait être rouge.
